\scanpage{31}
\mspage{29}
% archive sheet number: 30
$F_{\unsure{a|0}}$ et sur $F_b$. \inserted{C'est un épi, et} L'hom correspondant $\bar F\longrightarrow\overline{F'}$ est l'inclusion
\[
F_a\times_{F_{\unsure{0|*}}}F_b\longrightarrow F_a\times F_b
\]
qui n'est un épi que si $\alpha=\beta=0$.

\uline{Perplexité} : On devrait pouvoir définir une notion de dérivateur (additif) ``\uline{triangulé}'' (qui correspondrait à l'\unsure{intuition|intention} des cat. \deleted{\unsure{ab}} triangulées -- c'est un dérivateur satisfaisant à des axiomes assez forts), de sorte que pour les morphismes entre ceux-ci, l'exactitude à gauche et l'exactitude à droite soient \uline{\unsure{équivalentes|équivalents}}. \deleted{C'est peut-être}

\figure{\begin{tikzpicture}\draw[line width=2.8pt,line cap=round] (0,.38) -- (.23,0);\end{tikzpicture}}

\deleted{\unsure{Quest.}} \deleted{\unsure{Question de}}

``\uline{Théorèmes de changement de base}'' pour dérivateurs.

Soit $\mathbb D$ dérivateur à gauche (existence des $f_*$) et considérons
\[
f:I\longrightarrow J
\]
\inserted{dans $\mathrm{Diag}$}, soit $j_0\in\mathrm{Ob}\,J$, on se \deleted{\unsure{pr}}propose de ``calculer'' le foncteur
\[
F\longrightarrow f_*(F)_{j_0}=\alpha^*_{j_0}\big(f_*(F)\big)
\]
où
\[
\alpha_{j_0}:e\longrightarrow J
\]
est l'inclusion de la cat. ponctuelle $e$ dans $J$, de valeur \unsure{$j$.|$j_0$} \deleted{On voudrait trouver} \inserted{Considérons}
